\section{Discussion}
\label{sec:discussion}

\subsection{Economic Interpretation of the CVaR Reward}
\label{subsec:disc_cvar}

The superiority of the CVaR reward (A3, A4) over the Sharpe-based reward (A1, A2)
in the Brazilian equity market is consistent with two well-documented structural
features of emerging-market return distributions.
First, Brazilian equities exhibit pronounced negative skewness and excess kurtosis
relative to developed-market benchmarks, driven by periodic currency crises,
commodity price shocks, and fiscal uncertainty.
The Sharpe ratio, by treating upside and downside deviations symmetrically, implicitly
values a large positive outlier as much as it penalizes an equal negative one.
The CVaR reward, by contrast, focuses exclusively on the left tail and therefore
incentivizes the agent to avoid large drawdowns even at some cost in expected return.
In an environment where tail events are more frequent and severe than a normal
distribution would predict, this asymmetric objective yields portfolios that are
more appropriate for risk-averse investors.

Second, the parameter $\lambda = 0.5$ in the reward function
(Equation~\eqref{eq:cvar_reward}) creates a meaningful tension between return
maximization and tail-risk minimization.
The ablation results suggest that this balance is well calibrated: the CVaR-based
agents do not sacrifice too much expected return in pursuit of tail protection,
as evidenced by the competitive Calmar ratios relative to pure return-maximizing
strategies.

\subsection{Role of Macroeconomic State Variables}
\label{subsec:disc_macro}

The incremental performance of A4 over A3 confirms that macroeconomic
regime information is a valuable input to the portfolio allocation decision in the
Brazilian market.
The SELIC one-hot encoding provides the agent with a direct signal about the
monetary policy stance that is not fully captured by the trailing return and
volatility statistics already present in the state vector.
During rising-rate regimes, the combination of tighter financial conditions, reduced
earnings visibility, and elevated risk premiums creates an environment in which
active regime conditioning adds the most value.

The CDS proxy and credit spread play a complementary role, capturing dimensions
of credit-market stress that are correlated with but not identical to the interest-rate
regime.
Periods of elevated CDS spreads, for instance, can precede formal COPOM rate
decisions, giving the agent a forward-looking signal about macro deterioration
that the backward-looking SELIC one-hot encoding would miss until the next committee
meeting.

These findings are consistent with the theoretical argument that macro-financial
variables reduce the \emph{irreducible} component of portfolio uncertainty by
improving the precision of the agent's beliefs about the return-generating process.
From a practical standpoint, all three macro variables (SELIC direction, CDS proxy,
credit spread) are observable in real time, making the A4 agent directly implementable
without reliance on forecasts or latent state estimates.

\subsection{Limitations}
\label{subsec:disc_limitations}

Several limitations of this study should be acknowledged.

\textbf{Transaction costs.}
The evaluation protocol excludes transaction costs.
Given that the CQL agent rebalances quarterly---the same frequency as the
benchmarks---the additional cost from deviation of agent weights from benchmark
weights is unlikely to reverse the qualitative results, but it may reduce the
magnitude of outperformance, particularly for institutional investors subject to
market-impact costs on large B3 positions.
Future work should incorporate a realistic cost model, including bid-ask spreads
and market-impact functions.

\textbf{No short selling.}
The simplex constraint (Equation~\eqref{eq:action_space}) restricts the agent to
long-only positions.
This is consistent with regulatory constraints on many Brazilian retail and
institutional funds but prevents the agent from hedging systematic risk by taking
short positions in overvalued assets or index futures.
Extending the framework to a box-constrained action space would be a natural
next step.

\textbf{Regime shift risk.}
The regime classification is based on COPOM decisions, which are discrete events
occurring eight times per year.
In periods of rapid or unprecedented monetary-policy changes---such as the
emergency rate cuts during the COVID-19 pandemic---the regime label may lag
the true economic environment by several weeks.
More sophisticated regime-detection methods, such as hidden Markov models or
real-time macro nowcasting, could reduce this lag.

\textbf{Asset universe concentration.}
The ten-asset universe, while representative of the B3 blue-chip segment, is
highly concentrated in financials and commodity-linked sectors.
Expanding the universe to include small- and mid-cap stocks, fixed-income
instruments, or real-estate investment trusts (FIIs) would test the
generalizability of the proposed framework.

\subsection{Practical Implications}
\label{subsec:disc_implications}

For Brazilian fund managers, the results suggest that regime-aware portfolio
strategies offer meaningful improvements in downside risk management during
tightening cycles.
The largest gains from the A4 agent are concentrated precisely in the environment
most challenging for equity investors---rising rates---suggesting that the
incremental cost of computing and incorporating the macro state vector is well
justified by the reduction in tail risk during these periods.
The quarterly rebalancing frequency is compatible with the operational constraints
of most domestic equity funds and does not require intraday or high-frequency
monitoring.
The use of an offline RL framework furthermore implies that the strategy can be
backtested and validated on historical data without the methodological concerns
about look-ahead bias that complicate the evaluation of online RL agents, making
it more suitable for regulatory and audit scrutiny.
